% Part: computability
% Chapter: recursive-functions
% Section: examples

\documentclass[../../../include/open-logic-section]{subfiles}

\begin{document}

\olfileid{cmp}{rec}{exa}
\olsection{بنیادی بازگشتی فنکشناں دیاں مثالاں}


ساڈے کول بنیادی بازگشتی فنکشناں دیاں
کجھ مثالاں پہلے توں نیں: جمع تے ضرب
دے فنکشن~$\Add$ تے $\Mult$۔
شناختی فنکشن $\fn{id}(x) = x$ بنیادی
بازگشتی اے، کیونکہ ایہ بس~$\Proj{1}{0}$
اے۔ مستقل فنکشن $\fn{const}_n(x) = n$
وی بنیادی بازگشتی نیں، کیونکہ اوہناں نوں
$\Zero$ تے $\Succ$ توں لگاتار ترکیب
راہیں متعین کیتا جا سکدا اے۔ ایہ اوہدوں
کم آؤندا اے جدوں بنیادی بازگشتی تعریفاں
وچ مستقل عدد ورتنے ہوون۔ مثال لئی، جے
اسی فنکشن $f(x) = 2 \cdot x$
متعین کرنا چاہیئے، تاں اوہنوں
$\fn{const}_n(x)$ تے ضرب دی ترکیب توں
$f(x) =
\Mult(\fn{const}_2(x), \Proj{1}{0}(x))$
دے روپ وچ حاصل کر سکدے آں۔
اگوں اسی ایہ طریقہ ورتاں گے۔

\begin{prop}
  قوت چڑھاؤ دا فنکشن $\fn{exp}(x, y) = x^y$ بنیادی بازگشتی اے۔
\end{prop}

\begin{proof}
  اسی $\fn{exp}$ نوں بنیادی بازگشت
  راہیں ایویں متعین کر سکدے آں:
  \begin{align*}
    \fn{exp}(x, 0) & = 1\\
    \fn{exp}(x, y+1) & = \Mult(x, \fn{exp}(x,y)).
    \intertext{سختی نال آکھیے تاں ایہ بنیادی بازگشتی فنکشناں توں بازگشتی تعریف نہیں۔ باقاعدہ روپ ایہ اے:}
    \fn{exp}(x, 0) & = f(x)\\
    \fn{exp}(x, y+1) & = g(x, y, \fn{exp}(x,y)).
    \intertext{جتھے}
    f(x) & = \Succ(\Zero(x)) = 1\\
    g(x, y, z) & = \Mult(\Proj{3}{0}(x, y, z), \Proj{3}{2}(x, y, z)) = x \cdot z
  \end{align*}
  تے ایس لئی $f$ تے $g$ بنیادی
  بازگشتی فنکشناں توں ترکیب راہیں
  متعین ہوندے نیں۔
\end{proof}

\begin{prop}
  پیشرو فنکشن $\fn{pred}(y)$، جیہڑا
  ایویں متعین اے،
  \[
  \fn{pred}(y) = \begin{cases}
    0 & \text{جے $y=0$}\\
    y-1 & \text{ورنہ}
  \end{cases}
  \]
  بنیادی بازگشتی اے۔
\end{prop}

\begin{proof}
  غور کرو:
 \begin{align*}
   \fn{pred}(0) & = 0 \text{ تے}\\
   \fn{pred}(y+1) & = y.
 \end{align*}
 ایہ بنیادی بازگشتی تعریف دے لگ بھگ
 اے۔ سختی نال آکھیے تاں ایہ بنیادی
 بازگشت راہیں تعریف دے نمونے وچ پورا
 نہیں آؤندا، کیونکہ اوہ نمونہ گھٹّو گھٹ
 اک ہور دلیل~$x$ منگدا اے۔ ایہ گل وی
 وکھری اے کہ $\fn{pred}(y)$ نوں
 $\fn{pred}(y+1)$ دی تعریف وچ اصل وچ
 ورتیا ہی نہیں جاندا۔ پر اسی پہلے
 $\fn{pred}'(x, y)$ ایویں متعین کر
 سکدے آں:
 \begin{align*}
   \fn{pred}'(x, 0) & = \Zero(x) = 0,\\
   \fn{pred}'(x, y+1) & = \Proj{3}{1}(x, y, \fn{pred'}(x, y)) = y.
 \end{align*}
تے فیر اوہدے توں ترکیب راہیں
$\fn{pred}$ متعین کر سکدے آں،
مثلاً $\fn{pred}(x) = \fn{pred}'(\Zero(x), \Proj{1}{0}(x))$۔
\end{proof}

\begin{prop}
  عاملی فنکشن $\fn{fac}(x) = \fact{x} = 1 \cdot 2 \cdot 3
  \cdot \dots \cdot x$ بنیادی بازگشتی اے۔
\end{prop}

\begin{proof}
  بنیادی بازگشتی تعریف دا سیدھا
  روپ ایہ اے:
  \begin{align*}
    \fn{fac}(0) &= 1\\
    \fn{fac}(y+1) & = \fn{fac}(y) \cdot (y+1).
    \intertext{باقاعدہ طور اُتے سانوں پہلے اک دو-جگہ فنکشن $h$ متعین کرنا پیندا اے}
    h(x, 0) & = \fn{const}_1(x)\\
    h(x, y+1) & = g(x, y, h(x, y))
    \intertext{جتھے $g(x, y, z) = \Mult(\Proj{3}{2}(x, y, z),
      \Succ(\Proj{3}{1}(x, y, z)))$ تے فیر رکھو}
    \fn{fac}(y) & = h(\Proj{1}{0}(y), \Proj{1}{0}(y)) = h(y,y).
  \end{align*}
  ہن توں اگے اسی تھوڑی آسانی
  رکھاں گے تے ترکیب تے بنیادی
  بازگشت دیاں باقاعدہ تعریفاں
  ہر وار نہیں دیواں گے۔
\end{proof}

\begin{prop}
  صفر اُتے رک جان والی تفریق، $x \tsub y$،
  جیہڑی ایویں متعین اے،
  \[
  x \tsub y = \begin{cases}
    0 & \text{جے $x < y$}\\
    x-y & \text{ورنہ}
  \end{cases}
  \]
  بنیادی بازگشتی اے۔
\end{prop}

\begin{proof}
  ساڈے کول ایہ اے:
  \begin{align*}
    x \tsub 0 & = x\\
    x \tsub (y+1) & = \fn{pred}(x \tsub y) 
  \end{align*}
\end{proof}

\begin{prop}
 $x$ تے $y$ وچلا فاصلہ، $\left|x-y\right|$، بنیادی بازگشتی اے۔
\end{prop}

\begin{proof}
  ساڈے کول $\left| x-y \right| = (x \tsub y) + (y \tsub x)$
  اے؛ ایس لئی فاصلے نوں $+$ تے $\tsub$
  توں ترکیب راہیں متعین کیتا جا سکدا
  اے، تے ایہ دوہاں بنیادی بازگشتی نیں۔
\end{proof}

\begin{prop}
  $x$ تے $y$ وچوں وڈا عدد، $\fn{max}(x,y)$، بنیادی بازگشتی اے۔
\end{prop}

\begin{proof}
  اسی $\fn{max}(x,y)$ نوں $+$ تے
  $\tsub$ توں ترکیب راہیں ایویں
  متعین کر سکدے آں:
  \[
  \fn{max}(x,y) \defis x + (y \tsub x).
  \]
  جے $x$ وڈا اے، یعنی $x \ge y$،
  تاں $y \tsub x = 0$، ایس لئی $x
  + (y \tsub x) = x + 0 = x$۔
  جے $y$ وڈا اے، تاں $y \tsub x =
  y - x$، تے ایس لئی $x + (y \tsub x) = x + (y - x) = y$۔
\end{proof}

\begin{prop}
  \ollabel{prop:min-pr}
  $x$ تے $y$ وچوں چھوٹا عدد، $\fn{min}(x,y)$، بنیادی بازگشتی اے۔
\end{prop}

\begin{proof}
  مشق۔
\end{proof}

\begin{prob}
  \olref[cmp][rec][exa]{prop:min-pr} ثابت کرو۔
\end{prob}

\begin{prob}
دکھاؤ کہ \[
f(x, y) =
2^{(2^{\iddots^{2^{x}}})}\raisebox{1ex}{\bigg\rbrace}
\raisebox{1ex}{\text {$y$ وار $2$}}\] بنیادی بازگشتی اے۔
\end{prob}

\begin{prob}
دکھاؤ کہ صحیح عددی تقسیم $d(x, y) = \lfloor x/y \rfloor$
(یعنی تقسیم جتھے اعشاری نقطے توں اگے
دا سارا حصہ نظرانداز کیتا جائے) بنیادی
بازگشتی اے۔ جدوں $y = 0$ ہووے، اسی
$d(x, y) = 0$ مقرر کردے آں۔
بنیادی بازگشت تے ترکیب ورت کے~$d$
دی اک صاف تعریف دیو۔
\end{prob}

\begin{prop}
بنیادی بازگشتی فنکشناں دا سیٹ
ہیٹھلے دو عملاں تحت بند اے:
\begin{enumerate}
\item محدود جمع: جے $f(\vec x, z)$ بنیادی بازگشتی اے،
تاں ایہ فنکشن وی بنیادی بازگشتی اے:
\[
g(\vec x, y) \defis \sum_{z = 0}^y f(\vec x, z).
\]
\item محدود حاصلِ ضرب: جے $f(\vec x, z)$ بنیادی بازگشتی اے،
تاں ایہ فنکشن وی بنیادی بازگشتی اے:
\[
h(\vec x, y) \defis \prod_{z = 0}^y f(\vec x, z).
\]
\end{enumerate}
\end{prop}

\begin{proof}
مثال لئی، محدود جمع نوں ایہناں
مساواتاں نال بازگشتی طور اُتے
متعین کیتا جاندا اے:
\begin{align*}
  g(\vec x, 0) & = f(\vec x, 0)\\
  g(\vec x, y+1) & = g(\vec x, y) + f(\vec x, y+1).
\end{align*}
\end{proof}


\end{document}
